\begin{table}[htbp]
\centering
\caption{Robustness and Mechanism Tests}
\label{tab:robust}
\begin{tabular}{lccl}
\toprule
Specification & Coefficient & SE & Design \\
\midrule
\multicolumn{4}{l}{\textit{Panel A: Within-Sector Tests}} \\[3pt]
Female $\times$ Post (624 only, Log Emp.) & -0.0427*** & (0.0047) & DD \\
Female $\times$ Post (624 only, Earnings) & -40.8*** & (11.0) & DD \\
Female $\times$ Post (Manufacturing, Log Emp.) & 0.0426*** & (0.0045) & Placebo DD \\
\midrule
\multicolumn{4}{l}{\textit{Panel B: Sample Restrictions}} \\[3pt]
DDD, excl.\ COVID quarters & -0.1023*** & (0.0072) & DDD \\
DDD, high-allocation states & -0.0694*** & (0.0064) & DDD \\
DDD, low-allocation states & -0.1004*** & (0.0076) & DDD \\
\midrule
\multicolumn{4}{l}{\textit{Panel C: Grant Expiration}} \\[3pt]
DDD, active grant period & -0.0747*** & (0.0044) & DDD \\
DDD, post-expiration & -0.1067*** & (0.0079) & DDD \\
\bottomrule
\end{tabular}
\begin{tablenotes}[flushleft]
\small
\item \textit{Notes:} Panel A shows within-sector difference-in-differences (female vs.\ male) for Social Assistance (624) and a placebo test using Manufacturing (311, 332). Panel B restricts the sample: excluding COVID-affected quarters (2020Q2--2021Q3), and splitting by above/below median state ARP allocation per capita. Panel C decomposes the post period into active grants (2021Q4--2023Q3) and post-expiration (2023Q4--2024Q3). All specifications include cell, industry$\times$quarter, and state$\times$quarter fixed effects. Standard errors clustered at the state level. $^{***}p<0.01$, $^{**}p<0.05$, $^{*}p<0.1$.
\end{tablenotes}
\end{table}
